\section{Conclusion}
\label{sec:conclusion}

The analogy between confluence and flatness is exact once one compares the right things. Cycles in the reduction graph are the wrong things: they neither detect nor are detected by confluence. Confluence compares pairs of rewrite sequences from a common state; local confluence compares the two branches of each peak, and under termination the two conditions are equivalent. A peak becomes a closed boundary, and so something that can be filled, only when its branches have a common descendant. Recording peaks as generating branchings, and filling them when possible, turns the reduction graph into a reduction 2-complex on which elementary holonomy is defined. Elementary flatness is then exactly local confluence, and under termination it is exactly confluence. The proof of Newman's lemma reads as a tiling of a large region by elementary cells, and these statements are checked in Lean for arbitrary rewriting systems.

In string and term rewriting, critical pairs mark the branchings that can fail; disjoint and variable-overlap peaks are filled automatically. Working relative to chosen start terms, with concrete critical-pair instances and with verdicts certified only for completely explored reachable sets, the computations find exactly the expected obstructions on the curated examples. In two completion examples, the added rule leaves the branching in place and supplies its filler.

The central message can be stated without exaggeration. Local order-independence of reduction, the joinability of every peak, is exactly elementary flatness, and for terminating systems it is equivalent to full order-independence, that is, to confluence. On the reduction 2-complex, the local generators of nontrivial holonomy are the nonjoinable peaks, and in structured rewriting they are instances of critical pairs.
